% ===========================================================================
% Figure: Full system pipeline overview
% Data sources → ISensorSource → DataAligner → DepthFusion → GaussianModel → PLY
% ===========================================================================
\begin{figure*}[t]
\centering
\begin{tikzpicture}[node distance=0.8cm and 1.0cm]

% ---- Row 1: Sensor sources ----
\node[sensor] (lidar)  {LiDAR\\PointCloud};
\node[sensor, right=of lidar] (itof)   {iToF\\Depth};
\node[sensor, right=of itof]  (stereo) {Stereo\\RGB};
\node[sensor, right=of stereo] (imu)   {IMU};

% ---- Row 2: ISensorSource interface ----
\node[block, below=1.2cm of itof, minimum width=9cm, fill=blue!12]
  (iface) {\texttt{ISensorSource} (abstract interface)};

% ---- Row 3: Concrete implementations (side by side) ----
\node[blockgreen, below=0.8cm of iface, xshift=-2.8cm, minimum width=3.5cm]
  (reader) {DatasetReader\\(offline files)};
\node[blockorange, below=0.8cm of iface, xshift=2.8cm, minimum width=3.5cm]
  (orbbec) {OrbbecSource\\(live SDK)};

% ---- Row 4: Data aligner ----
\node[block, below=of iface, yshift=-2.2cm, minimum width=4cm]
  (aligner) {DataAligner\\(timestamp sync)};

% ---- Row 5: Depth fusion ----
\node[blockred, below=of aligner, minimum width=4cm]
  (fusion) {Depth Fusion\\(LiDAR + iToF + Stereo)};

% ---- Row 6: Depth completion ----
\node[blockorange, below=of fusion, minimum width=4cm]
  (completer) {SPNet Depth\\Completion (TensorRT)};

% ---- Row 7: Gaussian model ----
\node[blockgreen, below=of completer, minimum width=4cm]
  (gaussian) {GaussianModel\\(3DGS)};

% ---- Row 8: Output ----
\node[block, below=of gaussian, minimum width=4cm, fill=gray!10]
  (output) {PLY Map +\\Metrics};

% ---- Arrows: sensors → interface ----
\draw[arrow] (lidar.south) -- (iface.north -| lidar.south);
\draw[arrow] (itof.south)  -- (iface.north -| itof.south);
\draw[arrow] (stereo.south)-- (iface.north -| stereo.south);
\draw[arrow] (imu.south)   -- (iface.north -| imu.south);

% ---- Arrows: interface → implementations ----
\draw[arrow] (iface.south -| reader.north) -- (reader.north);
\draw[arrow] (iface.south -| orbbec.north) -- (orbbec.north);

% ---- Arrows: implementations → aligner ----
\draw[arrow] (reader.south) |- (aligner.west);
\draw[arrow] (orbbec.south) |- (aligner.east);

% ---- Main pipeline arrows ----
\draw[arrow] (aligner) -- (fusion);
\draw[arrow] (fusion)  -- (completer);
\draw[arrow] (completer) -- (gaussian);
\draw[arrow] (gaussian) -- (output);

% ---- Group: Pluggable layer ----
\begin{scope}[on background layer]
  \node[group=Pluggable Sensor Layer, fit=(iface)(reader)(orbbec),
        fill=blue!3] (pluggable) {};
\end{scope}

\end{tikzpicture}
\caption{Overview of the \winlic{} processing pipeline. Multi-modal sensor
data enters through the pluggable \texttt{ISensorSource} abstraction, is
synchronized by the DataAligner, fused into a dense depth map, completed by
the SPNet TensorRT engine, and fed into the 3D Gaussian Splatting model for
reconstruction.}
\label{fig:pipeline}
\end{figure*}
